\documentclass[12pt,a4paper]{article}

% ============================================================
% PACKAGES
% ============================================================
\usepackage[utf8]{inputenc}
\usepackage[T1]{fontenc}
\usepackage{lmodern}
\usepackage[french]{babel}
\usepackage{geometry}
\usepackage{graphicx}
\usepackage[table]{xcolor}
\usepackage{listings}
\usepackage{booktabs}
\usepackage{array}
\usepackage{fancyhdr}
\usepackage{titlesec}
\usepackage{enumitem}
\usepackage{mdframed}
\usepackage{float}
\usepackage{microtype}
\usepackage{setspace}
\usepackage{hyperref}

% ============================================================
% MISE EN PAGE
% ============================================================
\geometry{
  headheight=25pt,
  top=2.6cm,
  bottom=2.8cm,
  left=2.8cm,
  right=2.8cm
}
\setstretch{1.12}

% ============================================================
% COULEURS
% ============================================================
\definecolor{primary}{RGB}{28, 90, 168}
\definecolor{primarylight}{RGB}{88, 145, 220}
\definecolor{secondary}{RGB}{232, 240, 252}
\definecolor{codebg}{RGB}{248, 250, 255}
\definecolor{codeframe}{RGB}{170, 195, 230}
\definecolor{success}{RGB}{34, 130, 90}
\definecolor{warning}{RGB}{200, 120, 40}
\definecolor{darkgray}{RGB}{45, 55, 72}
\definecolor{lightgray}{RGB}{95, 110, 140}

% ============================================================
% STYLE DU CODE
% ============================================================
\lstset{
  backgroundcolor=\color{codebg},
  frame=single,
  rulecolor=\color{codeframe},
  basicstyle=\ttfamily\footnotesize,
  keywordstyle=\color{primary}\bfseries,
  commentstyle=\color{lightgray}\itshape,
  stringstyle=\color{success},
  breaklines=true,
  showstringspaces=false,
  numbers=left,
  numberstyle=\tiny\sffamily\textcolor{lightgray},
  numbersep=6pt,
  tabsize=2,
  captionpos=b,
  aboveskip=1em,
  belowskip=1em,
  xleftmargin=0pt,
  framexleftmargin=0pt
}

% ============================================================
% EN-TETE ET PIED DE PAGE
% ============================================================
\pagestyle{fancy}
\fancyhf{}
\fancyhead[L]{\small\sffamily\textcolor{primary!85!black}{Projet M1 GIL --- Gestion de flotte}}
\fancyhead[R]{\small\sffamily\textcolor{primary!85!black}{Univ. Rouen Normandie \textperiodcentered{} 2025--2026}}
\fancyfoot[C]{\small\sffamily\textcolor{primary!70!black}{\thepage}}
\renewcommand{\headrulewidth}{0.4pt}
\renewcommand{\headrule}{\hbox to\headwidth{\color{primary!35}\leaders\hrule height \headrulewidth\hfill}}
\renewcommand{\footrulewidth}{0pt}

% ============================================================
% STYLE DES TITRES
% ============================================================
\titleformat{\section}
  {\Large\sffamily\bfseries\color{primary}}
  {\thesection.}{0.55em}{}

\titleformat{\subsection}
  {\large\sffamily\bfseries\color{primary!85!black}}
  {\thesubsection.}{0.45em}{}

\titleformat{\subsubsection}
  {\normalsize\sffamily\bfseries\color{primary!75!black}}
  {\thesubsubsection.}{0.4em}{}

\titlespacing*{\section}{0pt}{1.6ex plus .2ex}{1ex plus .1ex}
\titlespacing*{\subsection}{0pt}{1.3ex plus .15ex}{0.7ex plus .1ex}

% ============================================================
% BOITE INFO
% ============================================================
\newmdenv[
  backgroundcolor=secondary,
  linecolor=primary,
  linewidth=1.1pt,
  leftline=true,
  rightline=false,
  topline=false,
  bottomline=false,
  innerleftmargin=11pt,
  innerrightmargin=11pt,
  innertopmargin=9pt,
  innerbottommargin=9pt,
  skipabove=0.6em,
  skipbelow=0.6em
]{infobox}

\hypersetup{
  colorlinks=true,
  linkcolor=primary!90!black,
  urlcolor=primarylight,
  pdfborderstyle={/S/U/W 0}
}

% ============================================================
% DOCUMENT
% ============================================================
\begin{document}

% ------------------------------------------------------------
% PAGE DE TITRE
% ------------------------------------------------------------
\begin{titlepage}
  \centering
  \vspace*{1cm}

  {\Large\textbf{Universit\'{e} de Rouen Normandie}\\[0.3cm]
  \large Master 1 -- G\'{e}nie Informatique et Logiciel}

  \vspace{1.5cm}
  {\color{primary!55}\rule{\linewidth}{0.5mm}}\\[0.45cm]

  {\Huge\sffamily\bfseries\color{primary} Projet Gestion de Flotte\\[0.35cm]}
  {\LARGE\sffamily\bfseries\color{primary!80!black} Rapport de Semaine 6}\\[0.35cm]
  {\large\sffamily\color{primary!65!black} Architecture Frontend \& S\'{e}curit\'{e} (SSO)}

  {\color{primary!55}\rule{\linewidth}{0.5mm}}\\[1.5cm]

  \begin{minipage}{0.45\textwidth}
    \begin{flushleft}
      \textbf{\'{E}tudiants :}\\
      Ahc\`{e}ne Amouchas \\
      Florian Pepin \\[0.5cm]
      \textbf{Encadrants :}\\
      Lydia \& Luc
    \end{flushleft}
  \end{minipage}
  \hfill
  \begin{minipage}{0.45\textwidth}
    \begin{flushright}
      \textbf{Ann\'{e}e universitaire :}\\
      2025 -- 2026 \\[0.5cm]
      \textbf{Semaine :}\\
      6 / 7
    \end{flushright}
  \end{minipage}

  \vfill
  {\large 16 avril 2026}
\end{titlepage}

\tableofcontents
\newpage

% ============================================================
\section{Introduction}
% ============================================================

La semaine 6 marque une \'{e}tape d\'{e}cisive dans le d\'{e}veloppement de la plateforme de gestion de flotte avec la mise en place de l'infrastructure frontend. L'objectif principal a \'{e}t\'{e} de concevoir une architecture capable de supporter la croissance du projet tout en garantissant une s\'{e}curit\'{e} robuste et centralis\'{e}e.

L'approche retenue s'appuie sur le concept de \textbf{Micro-Frontends} et l'utilisation d'un \textbf{Monorepo}, permettant une s\'{e}paration stricte des responsabilit\'{e}s et une scalabilit\'{e} organisationnelle optimale.

% ============================================================
\section{Semaine 6 : Architecture Frontend et S\'{e}curit\'{e} (SSO)}
% ============================================================

\subsection{Objectifs et Philosophie d'Architecture}

Pour r\'{e}pondre aux exigences de scalabilit\'{e} et de maintenabilit\'{e}, l'interface utilisateur n'a pas \'{e}t\'{e} con\c{c}ue comme une application monolithique, mais selon une architecture orient\'{e}e composants isol\'{e}s. Cette approche s'articule autour de deux concepts majeurs : le \textbf{Monorepo} et les \textbf{Micro-Frontends (MFE)}.

L'objectif est de permettre \`{a} plusieurs d\'{e}veloppeurs d'op\'{e}rer en parall\`{e}le sur des modules distincts (ex: cartographie, gestion des v\'{e}hicules) sans g\'{e}n\'{e}rer de conflits, tout en garantissant une coh\'{e}rence visuelle et s\'{e}curitaire centralis\'{e}e.

\subsection{Structure en Monorepo (NPM Workspaces)}

L'\'{e}cosyst\`{e}me frontend est regroup\'{e} au sein d'un d\'{e}p\^{o}t unique organis\'{e} via les \textbf{NPM Workspaces}. Cette structure permet de mutualiser les d\'{e}pendances globales tout en s\'{e}parant strictement le code source. L'arborescence se divise en deux cat\'{e}gories principales :

\begin{itemize}[leftmargin=*]
  \item \textbf{Le r\'{e}pertoire \texttt{apps/} :} Contient les applications ex\'{e}cutables de mani\`{e}re autonome (ex: \texttt{shell}, \texttt{location}, \texttt{vehicles}).
  \item \textbf{Le r\'{e}pertoire \texttt{packages/} :} Regroupe le code transversal partag\'{e} entre les applications (ex: \texttt{shared-ui} pour les composants graphiques, \texttt{shared-auth} pour la s\'{e}curit\'{e}).
\end{itemize}

\subsection{Approche Micro-Frontend via Module Federation}

L'assemblage des diff\'{e}rentes interfaces s'effectue dynamiquement c\^{o}t\'{e} client gr\^{a}ce au plugin \textbf{Module Federation} de l'outil de build Vite.js.

\begin{itemize}[leftmargin=*]
  \item \textbf{L'application H\^{o}te (Shell) :} Elle agit comme une coquille vide de haut niveau. Son r\^{o}le se limite \`{a} g\'{e}rer le routage global, l'interface racine (barre de navigation) et l'application stricte de la politique de s\'{e}curit\'{e}.
  \item \textbf{Les applications Distantes (Remotes) :} Ce sont des micro-applications ind\'{e}pendantes. Au moment de la navigation, le Shell t\'{e}l\'{e}charge et injecte dynamiquement (via \texttt{React.lazy} et \texttt{Suspense}) le composant expos\'{e} par le Remote.
\end{itemize}

\subsection{S\'{e}curit\'{e} Centralis\'{e}e et Single Sign-On (SSO) avec Keycloak}

Afin de garantir un niveau de s\'{e}curit\'{e} optimal, le processus d'authentification a \'{e}t\'{e} d\'{e}l\'{e}gu\'{e} au standard \textbf{OAuth2 / OpenID Connect (OIDC)} via un serveur \textbf{Keycloak} conteneuris\'{e}.

\begin{infobox}
\textbf{Flux d'authentification :} Aucun formulaire de connexion n'est impl\'{e}ment\'{e} dans l'application React. Les requ\^{e}tes non authentifi\'{e}es sont redirig\'{e}es vers Keycloak (Standard Flow avec PKCE). Apr\`{e}s validation, Keycloak restitue un jeton JWT permettant aux applications frontend d'habiliter l'acc\`{e}s en fonction des r\^{o}les (RBAC).
\end{infobox}

\subsection{Conception du Package Transversal : AuthProvider}

La m\'{e}canique d'interception et de v\'{e}rification du jeton OIDC a \'{e}t\'{e} encapsul\'{e}e dans le paquetage \texttt{@flotte/shared-auth}. Ce module expose un Provider de contexte React inject\'{e} \`{a} la racine de l'application Shell.

\begin{enumerate}
  \item \textbf{Verrou d'initialisation (Anti-Race Condition) :} Un verrou m\'{e}moris\'{e} hors de la port\'{e}e du composant garantit une unique instanciation du client Keycloak, m\^{e}me en environnement de d\'{e}veloppement (React Strict Mode).
  \item \textbf{Blocage du rendu :} Le Provider bloque le montage de l'arbre DOM tant que l'initialisation de Keycloak n'est pas r\'{e}solue, emp\^{e}chant l'affichage momentan\'{e} de routes prot\'{e}g\'{e}es (\textit{Flickering}).
  \item \textbf{Hook personnalis\'{e} (\texttt{useAuth}) :} Expose une API de fa\c{a}de permettant \`{a} n'importe quel micro-frontend de r\'{e}cup\'{e}rer l'identit\'{e} de l'utilisateur ou de d\'{e}clencher une d\'{e}connexion.
\end{enumerate}

\subsection{Routage Global et Chargement Asynchrone (Shell)}

L'application h\^{o}te (\texttt{Shell}) orchestre la navigation gr\^{a}ce \`{a} \texttt{react-router-dom}. Afin d'optimiser les performances, les modules distants ne sont pas charg\'{e}s au d\'{e}marrage mais import\'{e}s dynamiquement via \texttt{React.lazy()}. Ce m\'{e}canisme de \textit{Lazy Loading} diff\`{e}re le t\'{e}l\'{e}chargement des bundles JavaScript jusqu'au moment pr\'{e}cis o\`{u} l'utilisateur acc\`{e}de \`{a} la route correspondante.

\subsection{Communication Inter-Services : Client GraphQL Partag\'{e}}

Pour connecter l'interface \`{a} la Gateway GraphQL, un second paquetage partag\'{e} a \'{e}t\'{e} d\'{e}velopp\'{e} : \texttt{@flotte/shared-client}, bas\'{e} sur \texttt{@apollo/client}.

Gr\^{a}ce \`{a} un intercepteur r\'{e}seau (\textit{Context Link Apollo}), le client r\'{e}cup\`{e}re dynamiquement le jeton JWT fourni par l'\texttt{AuthProvider} et l'injecte dans l'en-t\^{e}te HTTP \texttt{Authorization: Bearer <token>} de chaque requ\^{e}te sortante.

\subsection{D\'{e}ploiement et Int\'{e}gration du Micro-Frontend ``V\'{e}hicules''}

Le d\'{e}veloppement du module de gestion de flotte (CRUD V\'{e}hicules) valide l'approche f\'{e}d\'{e}r\'{e}e. C'est un projet Vite.js autonome qui consomme l'API GraphQL pour afficher les donn\'{e}es m\'{e}tiers.

\begin{infobox}
\textbf{Configuration technique :}
\begin{itemize}
  \item \textbf{Allocation de ports :} Utilisation de ports TCP statiques (ex: 5002) pour \'{e}viter les conflits CORS.
  \item \textbf{Dependency Sharing :} Partage des biblioth\`{e}ques critiques (\texttt{react}, \texttt{apollo}) comme singletons pour \'{e}viter le double chargement et les conflits de Hooks.
\end{itemize}
\end{infobox}

\subsection{Commandes Fondamentales de l'Infrastructure}

La gestion du Monorepo NPM et le cycle de construction des Micro-Frontends s'articulent autour des commandes suivantes :

\begin{lstlisting}[language=bash, caption=Gestion du Monorepo et des MFEs]
# Creation d'une nouvelle sous-application
npm create vite@latest nom_app -- --template react-ts

# Installation du plugin de federation
npm install @originjs/vite-plugin-federation --save-dev

# Compilation isolee du module distant
npm run build --workspace=nom_app

# Exposition locale du module compile (preview)
npm run preview --workspace=nom_app

# Demarrage du Shell principal
npm run dev --workspace=shell
\end{lstlisting}

\subsection{Int\'{e}gration et Consommation du Contrat GraphQL}

La traduction du sch\'{e}ma GraphQL en composants React s'articule autour de quatre axes :
\begin{itemize}
  \item \textbf{Typage statique :} Les types GraphQL (ex: \texttt{Vehicle}) sont transpos\'{e}s en interfaces TypeScript.
  \item \textbf{Optimisation r\'{e}seau :} Utilisation de la syntaxe \texttt{gql} pour ne r\'{e}clamer que les champs n\'{e}cessaires (\textit{No Over-fetching}).
  \item \textbf{Rendu pilot\'{e} par les Enums :} Utilisation des \texttt{VehicleStatus} pour conditionner dynamiquement le rendu visuel.
  \item \textbf{Abstraction du cycle de vie (\texttt{useQuery}) :} Gestion native des \'{e}tats de chargement et de mise en cache.
\end{itemize}

\subsection{M\'{e}thodologie d'impl\'{e}mentation : Cas du module ``Conducteurs''}

L'ajout d'un nouveau domaine m\'{e}tier suit un processus standardis\'{e} :

\subsubsection{Initialisation et d\'{e}pendances}
\begin{lstlisting}[language=bash]
npm create vite@latest drivers -- --template react-ts
cd drivers
npm install @apollo/client graphql
npm install @originjs/vite-plugin-federation --save-dev
\end{lstlisting}

\subsubsection{Configuration de la F\'{e}d\'{e}ration}
\begin{lstlisting}[language=Java, caption=drivers/vite.config.ts (extraits)]
export default defineConfig({
  plugins: [
    react(),
    federation({
      name: 'driversapp',
      filename: 'remoteEntry.js',
      exposes: {
        './DriverList': './src/App.tsx',
      },
      shared: ['react', 'react-dom', '@apollo/client', 'graphql'],
    }),
  ],
  server: { port: 5003, strictPort: true }
});
\end{lstlisting}

\subsubsection{D\'{e}finition du Contrat GraphQL}
\begin{lstlisting}[language=Java, caption=drivers/src/queries.ts]
export const GET_DRIVERS = gql`
  query GetDrivers {
    drivers {
      items {
        id
        firstname
        lastname
        status
      }
    }
  }
`;
\end{lstlisting}

\subsubsection{Int\'{e}gration dans le Shell}
\begin{lstlisting}[language=Java, caption=shell/vite.config.ts]
remotes: {
  driversapp: 'http://localhost:5003/assets/remoteEntry.js',
}
\end{lstlisting}

\begin{lstlisting}[language=Java, caption=shell/App.tsx]
const DriverList = React.lazy(() => import('driversapp/DriverList'));

<Route path="drivers" element={
  <Suspense fallback={<div>Chargement du module...</div>}>
    <DriverList />
  </Suspense>
} />
\end{lstlisting}

% ============================================================
\section{Conclusion et Perspectives de Finalisation}
% ============================================================

Bien que l'ossature technique des Micro-Frontends et la s\'{e}curit\'{e} centralis\'{e}e soient d\'{e}sormais op\'{e}rationnelles, certains modules n\'{e}cessitent encore des ajustements finaux. 

\begin{infobox}
\textbf{Note sur le rendu :} La partie relative \`{a} la \textbf{cartographie interactive} ainsi que les derni\`{e}res optimisations de l'interface sont en cours de finalisation. Afin de fournir un livrable complet et un rapport de qualit\'{e} sup\'{e}rieure, l'ensemble du projet (code source et documentation finale) sera officiellement soumis ce \textbf{samedi}. 
\end{infobox}

Ce d\'{e}lai suppl\'{e}mentaire nous permet d'assurer une int\'{e}gration parfaite entre le service de localisation (backend) et son rendu visuel (frontend), tout en affinant la coh\'{e}rence globale de la plateforme.

\end{document}
